% part: intuitionistic-logic
% chapter: semantics
% section: topological-semantics

\documentclass[../../../include/open-logic-section]{subfiles}

\begin{document}

\olfileid[hi]{int}{sem}{top}

\olsection{सांस्थितिक अर्थविज्ञान}

अंतःप्रज्ञावादी तर्कशास्त्र का अर्थविज्ञान देने का एक अन्य ढंग संस्थिति की
गणितीय धारणा का उपयोग करना है।

\begin{defn}
  मान लें कि $X$ एक समुच्चय है। $X$ पर \emph{संस्थिति} ऐसा समुच्चय
  $\Top{O} \subseteq \Pow{X}$ है जो नीचे के गुणों को संतुष्ट करता है।
  ~$\Top{O}$ के !!{element} संस्थिति के \emph{विवृत समुच्चय} कहलाते हैं।
  समुच्चय $X$ और~$\Top{O}$ को मिलाकर \emph{सांस्थितिक समष्टि} कहते हैं।
  \begin{enumerate}
  \item रिक्त समुच्चय और पूरी समष्टि विवृत हैं:
    $\emptyset$, $X \in \Top{O}$।
  \item विवृत समुच्चय परिमित सर्वनिष्ठों के अधीन संवृत हैं: यदि
    $U$, $V \in \Top{O}$, तो $U \cap V \in \Top{O}$।
  \item विवृत समुच्चय स्वेच्छ संघों के अधीन संवृत हैं: यदि प्रत्येक
    $i \in I$ के लिए $U_i \in \Top{O}$, तो
    $\bigcup \Setabs{U_i}{i \in I} \in \Top{O}$।
  \end{enumerate}
\end{defn}

यदि विवृत समुच्चयों का संग्रह संदर्भ से समझा जा सके, तो हम किसी संस्थिति के
लिए केवल $X$ लिख सकते हैं। फिर भी ध्यान दें कि $X$ को विवृत समुच्चय प्रदान
किए जाने के बाद ही उसे संस्थिति कहा जा सकता है।

\begin{defn}
  अंतःप्रज्ञावादी प्रतिज्ञप्तिक तर्कशास्त्र का \emph{सांस्थितिक प्रतिरूप}
  एक त्रिक $\mModel{X} = \tuple{X, \Top{O}, V}$ है, जहाँ $\Top{O}$,
  ~$X$ पर कोई संस्थिति और $V$ ऐसा फलन है जो प्रत्येक प्रतिज्ञप्तिक चर को
  ~$\Top{O}$ का कोई विवृत समुच्चय देता है।

  कोई सांस्थितिक प्रतिरूप~$\mModel{X}$ दिए जाने पर हम $\Prop{X}{!A}$ को
  आगमनात्मक रूप से इस प्रकार परिभाषित कर सकते हैं:
  \begin{enumerate}
  \item $\Prop{X}{\lfalse} = \emptyset$
  \item $\Prop{X}{p} = V(p)$
  \item $\Prop{X}{!A \land !B} = \Prop{X}{!A} \cap \Prop{X}{!B}$
  \item $\Prop{X}{!A \lor !B} = \Prop{X}{!A} \cup \Prop{X}{!B}$
  \item $\Prop{X}{!A \lif !B} = \Interior{(X \setminus \Prop{X}{!A}) \cup
    \Prop{X}{!B}}$
  \end{enumerate}
  यहाँ $\Interior{V}$ वह फलन है जो किसी समुच्चय $V \subseteq X$ को उसके
  \emph{अभ्यंतर}, अर्थात उसमें समाविष्ट सभी विवृत समुच्चयों के संघ, में
  भेजता है। दूसरे शब्दों में,
  \[
  \Interior{V} = \bigcup \Setabs{U}{U \subseteq V \text{ और } U \in \Top{O}}.
  \]
\end{defn}

ध्यान दें कि किसी समुच्चय का अभ्यंतर सदैव विवृत होता है, क्योंकि वह विवृत
समुच्चयों का संघ है। अतः $\Prop{X}{!A}$ सदैव विवृत समुच्चय है।

यद्यपि सांस्थितिक अर्थविज्ञान अत्यंत अमूर्त है, फिर भी उसे प्रेरित करने वाले
ढंग से समझा जा सकता है। मान लें कि $X$ के !!{element}, अथवा ``बिंदु'', वे
बिंदु हैं जिन पर कथनों का मूल्यांकन किया जा सकता है। जिन सभी बिंदुओं पर
$!A$ सत्य है उनका समुच्चय,~$!A$ द्वारा व्यक्त प्रतिज्ञप्ति है। बिंदुओं का
हर समुच्चय संभाव्य प्रतिज्ञप्ति नहीं है; केवल~$\Top{O}$ के !!{element} ऐसे हैं।
प्रत्येक~$X$ के लिए $!A \Entails !B$ तभी जब प्रत्येक ऐसे बिंदु पर $!B$ सत्य
हो जहाँ~$!A$ सत्य है; अर्थात
$\Prop{X}{!A} \subseteq \Prop{X}{!B}$। असंगत कथन~$\lfalse$ कभी सत्य नहीं
होता, अतः $\Prop{X}{\lfalse} = \emptyset$।

अंतःप्रज्ञावादी रूप से मान्य नियमों के सत्य होने के लिए $!B \land !C$,
$!B \lor !C$ और $!B \lif !C$ द्वारा व्यक्त प्रतिज्ञप्तियों का~$!B$ और~$!C$
द्वारा व्यक्त प्रतिज्ञप्तियों से क्या संबंध होना चाहिए; अर्थात
$!A \Proves !B$ तभी जब $\Prop{X}{!A} \subset \Prop{X}{!B}$? प्रत्येक $!A$
के लिए हमें $\lfalse \Proves !A$ चाहिए, जो इसलिए संतुष्ट है कि प्रत्येक
$U$ के लिए $\emptyset \subseteq U$। चूँकि $!B \land !C \Proves !B$, हमें
$\Prop{X}{!B \land !C} \subseteq \Prop{X}{!B}$ चाहिए, और इसी तरह
$\Prop{X}{!B \land !C} \subseteq \Prop{X}{!C}$। $W \subseteq U$ और
$W \subseteq V$ को संतुष्ट करने वाला सबसे बड़ा समुच्चय $U \cap V$ है।
विलोमतः $!B \Proves !B \lor !C$ और $!C \Proves !B \lor !C$, अतः हमें
$\Prop{X}{!B} \subseteq \Prop{X}{!B \lor !C}$ और
$\Prop{X}{!C} \subseteq \Prop{X}{!B \lor !C}$ चाहिए। ऐसा सबसे छोटा
समुच्चय~$W$ जिसमें $U \subseteq W$ और $V \subseteq W$, $U \cup V$ है।

$\lif$ की परिभाषा कठिन है: $!A \lif !B$ उस सबसे दुर्बल प्रतिज्ञप्ति को
व्यक्त करता है जो~$!A$ के साथ मिलकर~$!B$ को निष्पन्न करती है।
$(!A \lif !B) \land !A \Proves !B$ से स्पष्ट है कि $!A \lif !B$,~$!A$ के
साथ मिलकर~$!B$ को निष्पन्न करता है। अतः $\Prop{X}{!A \lif !B}$ ऐसा सबसे
बड़ा विवृत समुच्चय होना चाहिए कि
$\Prop{X}{!A \lif !B} \cap \Prop{X}{!A} \subset \Prop{X}{!B}$; इससे हमारी
परिभाषा मिलती है।

\end{document}
